% BC
The Betweenness Centrality benchmark is taken from the SSCA2 (Scalable
Synthetic Compact Application 2) v2.2 benchmark \cite{SSCA2};
specifically, we implement the ``fourth'' kernel in this benchmark,
the computation of {\em betweenness centrality}:

\begin{quotation}
The intent of this kernel is to identify the set of vertices in the
graph with the highest betweenness centrality score. Betweenness
Centrality is a shortest paths enumeration-based centrality
metric, introduced by Freeman (1977). This is done using a betweenness
centrality algorithm that computes this metric for every vertex in the
graph. Let $\sigma_{st}$ denote the number of shortest paths between
vertices $s$ and $t$, and $\sigma_{st}(v)$ the number of those paths
passing through $v$. Betweenness Centrality of a vertex $v$ is defined
as $BC(v) = \Sigma_{s\not=v\not=t \in V} \sigma_{st}(v)/\sigma_{st}$.

The output of this kernel is a betweenness centrality score for each
vertex in the graph and the set of vertices the highest betweenness
centrality score.
\end{quotation}

We implement the Brandes' algorithm described in the benchmark,
augmenting Dijkstra's single-source shortest paths (SSSP) algorithm,
for unweighted graphs. We have implemented the exact variant of the benchmark
(\code{K4approx=SCALE}). 

The solution we implement makes one very strong assumption: that the
graph is ``small'' enough to fit in the memory of a single
place. Since the graph itself is not modified during the execution of
this benchmark, it thus becomes possible to implement this benchmark
by replicating the graph across all places. Now effectively the
computation can be performed in an embarrassingly parallel fashion. The
set of $N$ vertices is statically partitioned among $P$ places; each place is
responsible for performing the computation for the {\em source}
vertices assigned to it (for all $N$ target vertices) and computes its
local \code{betweennessMap}. After this is
done an \code{allReduce} collective performs an \code{AND} summation
    across all local \code{betweennessMaps}. 

Clearly, statically partitioning the work amongst all places in this
way is possible but not ideal. The amount of work associated with one source node
$v$ vs another $w$ could be dramatically different. Consider for
instance a degenerate case: a graph of $N$ vertices, labeled $1\ldots
N$, with an edge $(i,j)$ if $i < j$. Clearly the work associated with
vertex $1$ is much more than the work associated with vertex $N$. 

We next discuss how to dynamically load-balance these
tasks across all places using the GLB framework. 

%% \subsubsection{Structure of the code}
%% The multi-place structure of this code is quite straightforward. We
%% use the standard at-scale X10 idiom for SPMD code:
%% \begin{enumerate}
%% \item Create a distributed data structure, accessed through a
%%   \code{PlaceLocalHandle} which points to an instance of a ``host''
%%   object (an instance of \code{SSCA2}) at each node. (Invocation of
%%   \code{makeFlat}.) 
%% \item Run the actual computation (invocation of \code{broadcastFlat}
%%   with the closure specifying the work to be done at each place.

%%   The work done in each place is the computation of a
%%   \code{bfsShortestPaths} for a slice of vertices from
%%   \code{startVertex} to \code{endVertex}. This loops through the
%%   vertices in the given range, computing the betweennessMap
%%   locally. 

%%   Note this computation is single threaded. It is not difficult to
%%   make this multi-threaded. 

%% \item Once the local computation has been done, and the local
%%   contribution to the betweennessMap computed, a
%%   \code{Team.WORLD.allreduce} ensures that the results are reduced and
%%   made available at place $0$. 
%% \end{enumerate}

%% \subsection{Validation.}
%% The benchmark calls for validating the algorithms by running it on a
%% graph (a torus)  with known analytical results for the betweenness
%% centrality.  Because of lack of time, we have not been able to do
%% this. If our submission is selected for presentation, we plan to
%% validate the algorithm by the time of the presentation.


%% \subsection{Performance measurement.}
%% We discuss the performance together with BC-G in Section~\ref{sec:bcg}.

